% Verbatim from /Users/alingalatan/Downloads/cubic_chaos_partial_results (2).zip:binary_cubic_proof.tex (16 November 2026); GPT (a ChatGPT session of the user), relayed by the user.
% Intended filepath: binary_cubic_proof.tex
\documentclass[11pt]{article}
\usepackage[margin=1in]{geometry}
\usepackage{amsmath,amssymb,amsthm}
\newtheorem{theorem}{Theorem}
\newtheorem{lemma}{Lemma}
\newcommand{\E}{\mathbb E}
\newcommand{\He}{\operatorname{He}}
\title{A uniform strict Gaussian $L^1/L^2$ bound for every binary cubic chaos}
\author{}
\date{}
\begin{document}
\maketitle
\noindent\textbf{Status.} Every step below is rigorous and elementary. The result
covers all third-chaos polynomials in two independent real Gaussian variables,
including arbitrary mixtures of all four cubic Hermite basis elements. It is
not a proof in arbitrary dimension.

\begin{theorem}\label{thm:binary}
For every real $h\in H_3(\mathbb R^2)$,
\[
 \boxed{\E|h|\leq\sqrt{\frac{35}{18\pi}}\,\|h\|_2
       =\sqrt{\frac2\pi}\frac{\sqrt{35}}6\,\|h\|_2.}
\]
In particular this is strictly stronger than the proposed Gaussian constant
for every nonzero binary cubic.
\end{theorem}

\begin{lemma}[An angular inequality]\label{lem:angular}
Let $\Theta$ be uniform on $[0,2\pi]$ and let
\[
 P(\theta)=a\cos\theta+b\sin\theta+c\cos3\theta+d\sin3\theta.
\]
Then
\[
 \E|P(\Theta)|\leq\frac{\sqrt{80}}{3\pi}
             \bigl(\E P(\Theta)^2\bigr)^{1/2}.
\]
The constant in this angular inequality is sharp.
\end{lemma}

\begin{proof}
The assertion is trivial when $P=0$. Otherwise put $f=\operatorname{sign}P$.
The trigonometric polynomial $P$ has degree at most three and
$P(\theta+\pi)=-P(\theta)$, so $f$ has either one or three sign changes
per half-circle. Tangential zeros cause no sign change. Degenerate cases
are also obtained by limits of polynomials with simple zeros.

For $k\in\mathbb Z$ put
$\widehat f(k)=\E[f(\Theta)e^{-ik\Theta}]$. Since $P$ only uses frequencies
$\pm1,\pm3$, Cauchy--Schwarz gives
\[
 \E|P|=\E fP
 \leq\bigl(2|\widehat f(1)|^2+2|\widehat f(3)|^2\bigr)^{1/2}
       (\E P^2)^{1/2}.
\]
We show that the parenthesis is at most $80/(9\pi^2)$.

If there is one sign change per half-circle, $f$ is a rotated signed
half-circle indicator, and $|\widehat f(k)|=2/(\pi k)$ for positive odd $k$.
This gives exactly $80/(9\pi^2)$.

Otherwise rotate a sign-change point to zero and write the other two in
the same half-circle as $u$ and $u+v$, where $u,v>0$ and $u+v<\pi$.
After multiplying $f$ by $-1$ if necessary, it is $+1,-1,+1$ on the three
successive intervals between these points and $\pi$. Direct integration,
or integration by parts using its jumps, gives for odd positive $k$
\[
 |\widehat f(k)|=\frac2{\pi k}
 \left|1-e^{iku}+e^{ik(u+v)}\right|.
\]
Define
\[
 S=\sin(u/2)\sin(v/2)\cos((u+v)/2)>0,
 \quad
 S_3=\sin(3u/2)\sin(3v/2)\cos(3(u+v)/2).
\]
Expanding the squared modulus gives
\[
 |1-e^{iu}+e^{i(u+v)}|^2=1-8S,
 \quad
 |1-e^{3iu}+e^{3i(u+v)}|^2=1-8S_3.
\]
We claim $S_3\geq-8S$. By triple-angle identities,
\[
 \frac{S_3}{S}
 =(1+2\cos u)(1+2\cos v)(2\cos(u+v)-1).
\]
Set $C=\cos((u+v)/2)$ and $D=\cos((u-v)/2)$, so $0<C\leq D\leq1$.
The right-hand side becomes
\[
 (4C^2-3)(4C^2+4D^2+4CD-3).
\]
If $4C^2-3\geq0$, the other factor is also positive, since $D\geq C$.
If $4C^2-3<0$, the second factor is increasing in $D$, so this product
is bounded below by its value at $D=1$, namely
$(4C^2-3)(1+2C)^2$. Finally,
\[
 (4C^2-3)(1+2C)^2+8
 =(2C-1)^2(4C^2+8C+5)\geq0.
\]
Thus $S_3\geq-8S$ as claimed. It follows that
\begin{align*}
 2|\widehat f(1)|^2+2|\widehat f(3)|^2
 &=\frac8{\pi^2}\left(\frac{10}{9}
                 -8\left(S+\frac{S_3}{9}\right)\right)\\
 &\leq\frac8{\pi^2}\left(\frac{10}{9}-\frac89 S\right)
 <\frac{80}{9\pi^2}.
\end{align*}
This proves the inequality. Sharpness follows by taking
$P(\theta)=\cos\theta-\frac13\cos3\theta$, whose sign is the sign of
$\cos\theta$; it is proportional to the projection of that sign function
onto frequencies $\pm1,\pm3$.
\end{proof}

\begin{proof}[Proof of Theorem~\ref{thm:binary}]
Write $(X,Y)=(R\cos\Theta,R\sin\Theta)$, where $R\sim\chi_2$ is independent
of uniform $\Theta$. Every homogeneous cubic decomposes uniquely into a
harmonic cubic plus $r^2$ times a linear form. Wick subtraction replaces
$r^2\ell$ by $(r^2-4)\ell$, since $\Delta(r^2\ell)=8\ell$ in dimension two.
Consequently every $h\in H_3(\mathbb R^2)$ has the form
\[
 h(R,\Theta)=A R^3\cos(3\Theta-\alpha)
       +B R(R^2-4)\cos(\Theta-\beta),
\]
with $A,B\geq0$ and arbitrary phases. For fixed $R$, orthogonality of
frequencies one and three gives
\[
 q(R):=\E_\Theta h(R,\Theta)^2
 =\frac12\left(A^2R^6+B^2R^2(R^2-4)^2\right).
\]
Lemma~\ref{lem:angular}, followed by weighted Cauchy--Schwarz, yields
\[
 \E|h|\leq\frac{\sqrt{80}}{3\pi}\E_R\sqrt{q(R)}
 \leq\frac{\sqrt{80}}{3\pi}
       \sqrt{\E R\;\E(q(R)/R)}.
\]
The chi-two moments satisfy
\[
 \E R=\sqrt{\pi/2},\qquad
 \E R^3=3\E R,\qquad \E R^5=15\E R,
\]
so
\[
 \E R\;\E(q(R)/R)=\frac\pi4(15A^2+7B^2).
\]
Also $\E R^2=2$, $\E R^4=8$, and $\E R^6=48$, hence
\[
 \|h\|_2^2=24A^2+8B^2.
\]
Therefore
\begin{align*}
 (\E|h|)^2
 &\leq\frac{80}{9\pi^2}\frac\pi4(15A^2+7B^2)\\
 &\leq\frac{80}{9\pi^2}\frac{7\pi}{32}(24A^2+8B^2)
 =\frac{35}{18\pi}\|h\|_2^2,
\end{align*}
as desired.
\end{proof}

\paragraph{Checks.}
The estimate includes $\He_3(X)$ by taking a polynomial independent of $Y$.
It includes every combination of
$\He_3(X)$, $(X^2-1)Y$, $X(Y^2-1)$, and $\He_3(Y)$; there is no reflection
or relative-phase restriction. No numerical quadrature or external
extremality assertion is used in the proof. The decimal value of the
constant, for orientation only, is $0.786724786568858\ldots$, less than
$\sqrt{2/\pi}=0.797884560802865\ldots$.
\end{document}
